% Werke X/1 (1917), native163–166, printed156–159.
\begingroup
\setlength{\headheight}{25pt}
\SourceLeaf{163}\ArithmeticHead{156}{L}{ANALYSIS.\quad NACHLASS.}
\[\frac{16}{\Pi^4}(1{,}014678)+\Sigma\frac{m^4-6m^2n^2+n^4}{(mm+nn)^4}\]
\[\frac{1-\frac12xx-\frac1{24}x^4-\frac1{240}x^6\ldots}{1+\frac12xx-\frac1{24}x^4+\frac1{240}x^6\ldots}
\left[=\cos\text{ lemn }x=\frac{\mu x}{\nu x}\right].\]
\NumberedItem{7}
\noindent{\small [S. 66]}
\par\indent\textls[120]{Zähler} bei \(\sin90^0\)\quad
\(\left[90^0\text{ heisst }\dfrac{360^0}{4},\text{ also }=\int_0^1\dfrac{dp}{\surd(1-p^4)}=\Pi\right]\)
\[\renewcommand{\arraystretch}{1.65}\begin{array}{rrl}
&1{,}31102878&[=\Pi]\\
-&6455202&\left[=\frac1{60}\Pi^5\right]\\
-&113514&\left[=\frac1{10080}\Pi^9\right]\\\cline{1-2}
&1{,}24534162&\\
+&300&\\\cline{1-2}
&1{,}24534462&[=M\Pi]
\end{array}\]
\noindent\textls[120]{Nenner}
\[\renewcommand{\arraystretch}{1.65}\begin{array}{rrl}
&1{,}2461884&\left[=1+\frac{\Pi^4}{12}\right]\\
-&86583&\left[=\frac1{10080}\Pi^8\right]\\\cline{1-2}
&1{,}24532257&\\
+&2196&\\\cline{1-2}
&1{,}24534453&[=N\Pi]
\end{array}\]
\[\begin{array}{cccl}
\left(N\dfrac{1[\varphi]}{1+\surd-1}\right)^2&-&\left(M\dfrac{1[\varphi]}{1+\surd-1}\right)^2&=\mu[\varphi]\\
&+&&=\nu[\varphi]{}^{[*)]}
\end{array}\]
\renewcommand{\thefootnote}{[*)}\footnotetext{In der Handschrift werden hier die Buchstaben \(m,n\) benutzt; wir haben in Übereinstimmung mit art. [9.] die Zeichen \(\mu,\nu\) gesetzt.]}
\[\left(N\frac{1[\varphi]}{1+\surd-1}\right)^4-\left(M\frac{1[\varphi]}{1+\surd-1}\right)^4=N(1+\surd-1)[\varphi].\]
\clearpage
\SourceLeaf{164}\fancyhead{}\fancyhead[R]{\small157}\fancyhead[C]{\scriptsize\shortstack{ÄLTESTE UNTERSUCHUNGEN ÜBER\\ LEMNISKATISCHE FUNKTIONEN.\quad II.}}
\noindent Hinc
\[N^4+M^4=N2[\varphi].\]
\NumberedItem{8}
\noindent{\small S. 67]}
\par\indent Sit \(1{,}24\ldots=a[=M\Pi=N\Pi]\)
\[\renewcommand{\arraystretch}{1.3}\begin{gathered}
N\varpi=2a^4{}^{[*)]}\\
N2\varpi=(2a^4)^4\\
N4\varpi=(2a^4)^{16}\\
N8\varpi=(2a^4)^{64}\\
\text{etc.}\\
N\lambda\varpi=(2a^4)^{\lambda\lambda}\\[6pt]
M2[\varphi]=2MN\surd(N^4-M^4)
\end{gathered}\]
\renewcommand{\thefootnote}{[*)}\footnotetext{In der Handschrift steht in den folgenden Formeln statt \(\varpi\) das gewöhnliche \(\Pi\), so daß also dieser Buchstabe hier das Doppelte der Größe bedeuten würde, die er oben art. [3.], S. 64 des \textsc{Leiste}, bezeichnet; es wurde darum hier die später von \textsc{Gauss} ständig benutzte Bezeichnung
\[\varpi=2\int_0^1\frac{dp}{\sqrt{1-p^4}}\]
angewendet.]}
\noindent{\small [S. 70]}
\[\renewcommand{\arraystretch}{1.7}\begin{array}{rl}
[\nu^2=]MM+NN={}&1+1[\varphi^2]+\frac16[\varphi^4]-\frac1{30}[\varphi^6]+\frac{17}{2520}[\varphi^8]\ldots\\\hline
\nu={}&1+\frac12[\varphi^2]-\frac1{24}[\varphi^4]+\frac1{240}[\varphi^6]+\frac{17}{40320}[\varphi^8]\ldots
\end{array}\]
\[\begin{gathered}
N2[\varphi]=M^4+N^4\\
M2[\varphi]=2MN\mu\nu\\[5pt]
\mu\nu=N(1+\surd-1)[\varphi]\\[5pt]
\mu\mu=\frac{NN-MM}{NN+MM}\nu\nu
\end{gathered}\]
\[\begin{aligned}
\mu2[\varphi].\nu2[\varphi]&=M(1+\surd-1)[\varphi]^4+N(1+\surd-1)[\varphi]^4\\
&[=N2(1+\surd-1)\varphi].
\end{aligned}\]
\clearpage
\SourceLeaf{165}\ArithmeticHead{158}{L}{ANALYSIS.\quad NACHLASS.}
\noindent{\small [S. 71]}
\[\begin{aligned}
N4[\varphi]={}&M^{16}+4M^{12}N^4+6M^8N^8+4M^4N^{12}+N^{16}\\
&+16(N^8\nu^8-N^8\nu^6\mu^2-\nu^8N^6M^2).
\end{aligned}\]
\NumberedItem{9}
\noindent{\small [S. 71]}
\par\indent Bei \(45^0\left[=\frac{\varpi}{4}\right]\)\par
\noindent\begin{minipage}[t]{.52\linewidth}
\begin{center}Zähler des Sinus\end{center}
\[\renewcommand{\arraystretch}{1.6}\begin{array}{rrl}
&0{,}65551439&\left[=\frac{\varpi}{4}\right]\\
-&201725&\\-&221&\\\cline{1-2}
&0{,}65349493&\left[=M\frac{\varpi}{4}\right]
\end{array}\]
\end{minipage}\hfill\begin{minipage}[t]{.43\linewidth}
\begin{center}Nenner des Sinus\end{center}
\[\renewcommand{\arraystretch}{1.6}\begin{array}{rrl}
&1{,}01538678&\\-&338&\\\cline{1-2}
&1{,}01538340&\left[=N\frac{\varpi}{4}\right]
\end{array}\]
\end{minipage}
{\small\[\begin{array}{rl@{\qquad}rl}
\left[\log_{10}M\frac{\varpi}{4}=\right]&9{,}8152422&0{,}0066301&\left[=\log_{10}N\frac{\varpi}{4}\right]\\[10pt]
\left[4\log_{10}M\frac{\varpi}{4}=\right]&2609688&0{,}0265204&\left[=4\log_{10}N\frac{\varpi}{4}\right]
\end{array}\]}
\[\renewcommand{\arraystretch}{1.6}\begin{array}{rl}
\left[\left(M\frac{\varpi}{4}\right)^4=0{,}\right]&18237647\\
\left[\left(N\frac{\varpi}{4}\right)^4=\right]&1{,}06291846\\\cline{2-2}
\left[M\frac{\varpi}{2}=\right]&1{,}24534493,
\end{array}\]
\noindent welches mithin die \textls[120]{erste} Rechnung oben [siehe art. 3] bestätigt.
\[\Pi=1{,}245344\left[=M\frac{\varpi}{2}=N\frac{\varpi}{2}=a,\text{ siehe art. 8}\right]\]
\[\begin{array}{rl@{\qquad}rl}
l\Pi&0{,}0952896&&\\[8pt]
[4\log_{10}\Pi=]&0{,}3811584&2{,}40524&[=\Pi^4]\\[8pt]
&&4{,}81048&=N\varpi\ [=2\Pi^4]
\end{array}\]
\[\log\text{ hyp. dieser Zahl }=1{,}5708=\tfrac12\pi\text{ Circuli?}\]
\clearpage
\SourceLeaf{166}\fancyhead{}\fancyhead[R]{\small159}\fancyhead[C]{\scriptsize\shortstack{ÄLTESTE UNTERSUCHUNGEN ÜBER\\ LEMNISKATISCHE FUNKTIONEN.\quad II.}}
\NumberedItem{10}
\noindent{\small [S. 72--73]}
{\small\[\renewcommand{\arraystretch}{1.75}\begin{array}{rrrrl}
\log\text{ Denom[inatoris]}={}&\frac1{12}x^4&-\frac1{10080}x^8&+\frac{17}{19958400}x^{12}&\ldots\\
&&-\frac1{288}x^8&+\frac1{120960}[x^{12}]&\ldots\\
&&&+\frac1{5184}[x^{12}]&\ldots\\\cline{2-5}
[\log N=]&\frac1{12}x^4&-\frac1{280}x^8&+\frac1{4950}x^{12}&-\frac1{78000}x^{16}\ldots
\end{array}\]}
\[\log\text{ Numeratoris}=\log x-\frac1{60}x^4-\frac1{4200}x^8-\frac1{321750}x^{12}-\frac1{19890000}x^{16}\ldots\]
\[d\log\sin=\frac{\surd(1-\sin^4)}{\sin}\,d\text{ arc}\]
\[s=\frac{\varphi}{1+\frac1{10}s^4\ldots}\]
\[s=\varphi-\left(\frac1{10}s^5+\frac1{24}s^9\ldots\right)\]
\[ls=l\varphi-\frac1{10}[\varphi^4]+\frac1{300}[\varphi^8]-\frac1{4875}[\varphi^{12}]\ldots\]
\ClosingRule
\clearpage
\endgroup
